% 数学ノート集・共通テンプレート v2（積分論2026年度第1部の版面を参考）
% A5判。章内で定義・定理・例・注意等の番号を共通化。
% コンパイル： platex -interaction=nonstopmode math_note_template_v2.tex (2回)
%              dvipdfmx math_note_template_v2.dvi
\documentclass[a5paper,10pt,onecolumn]{jsarticle}
\usepackage[dvipdfmx,top=22mm,bottom=20mm,inner=19mm,outer=18mm,headheight=13pt,headsep=6mm]{geometry}
\usepackage{amsmath,amssymb,amsthm,mathtools}
\usepackage{enumitem,fancyhdr,titlesec}
\usepackage[dvipdfmx,bookmarks=true,colorlinks=false]{hyperref}
\usepackage{pxjahyper}

% ----- 差し替える文書情報 -----
\newcommand{\NoteTitle}{Brown運動の構成と基本性質}
\newcommand{\NoteSeries}{数学ノート集}

% ----- 文字・組版 -----
% 参考PDFの質素な明朝＋Computer Modernの印象を引き継ぎながら，
% A5判で窮屈にならないよう本文の行間を少し確保する．
\linespread{1.13}
\setlength{\parindent}{1zw}
\setlength{\parskip}{0pt}
\setlength{\abovedisplayskip}{6pt plus 2pt minus 2pt}
\setlength{\belowdisplayskip}{6pt plus 2pt minus 2pt}
\setlength{\abovedisplayshortskip}{3pt plus 2pt minus 1pt}
\setlength{\belowdisplayshortskip}{4pt plus 2pt minus 1pt}
\setlist[enumerate]{topsep=3pt,itemsep=1pt,parsep=0pt,leftmargin=2.4zw}
\titleformat{\section}{\normalfont\normalsize\bfseries}{\thesection}{0.8zw}{}
\titlespacing*{\section}{0pt}{2.0ex plus .4ex}{.9ex}
\titleformat{\subsection}{\normalfont\normalsize\bfseries}{\thesubsection}{0.7zw}{}
\titlespacing*{\subsection}{0pt}{1.5ex plus .3ex}{.6ex}
\raggedbottom

% ----- 柱とページ番号：ページ上部、罫線なし -----
\fancypagestyle{mathnote}{%
  \fancyhf{}%
  \fancyhead[L]{\footnotesize \NoteSeries\ （\NoteTitle）}%
  \fancyhead[R]{\footnotesize \thepage}%
  \renewcommand{\headrulewidth}{0pt}%
  \renewcommand{\footrulewidth}{0pt}%
}
\pagestyle{mathnote}

% ----- 節内通し番号の定義・定理群 -----
\newtheoremstyle{mathnotetheorem}
 {4.5pt}{4.5pt}{\normalfont}{}%
 {\bfseries}{}{0.65zw}%
 {\thmname{#1}\thmnumber{ #2}\thmnote{（#3）}}
\theoremstyle{mathnotetheorem}
\newtheorem{theorem}{定理}[section]
\newtheorem{definition}[theorem]{定義}
\newtheorem{lemma}[theorem]{補題}
\newtheorem{proposition}[theorem]{命題}
\newtheorem{corollary}[theorem]{系}
\newtheorem{example}[theorem]{例}
\newtheorem{remark}[theorem]{注意}
\renewcommand{\proofname}{証明}
\renewcommand{\qedsymbol}{$\square$}
\makeatletter
\renewenvironment{proof}[1][\proofname]{%
  \par\pushQED{\qed}\normalfont\topsep4\p@\@plus2\p@\relax
  \trivlist\item[\hskip\labelsep\bfseries #1]\ignorespaces
}{\popQED\endtrivlist\@endpefalse}
\makeatother

% ----- 使うノートで適宜拡張する記号 -----
\newcommand{\R}{\mathbb{R}}
\newcommand{\norm}[1]{\lVert #1\rVert}
\newcommand{\abs}[1]{\lvert #1\rvert}

\begin{document}
\section*{読了後にできること}
本稿を読むための前提は，Brown運動の独立増分と正規分布．
本稿の具体的な到達課題は，Brown運動の共分散と二乗平均増分を直接求める．

\section{確率過程とGaussian系}

$(\Omega,\mathcal F,P)$ を確率空間とする．
本稿では $P(A)$ を事象 $A$ の確率，
$E[Z]=\int_\Omega Z\,dP$ を可積分確率変数 $Z$ の期待値，
$\operatorname{Cov}(X,Y)=E[(X-E[X])(Y-E[Y])]$ を共分散と書く．
また $\langle u,v\rangle$ は有限次元Euclid空間の標準内積とする．
「ほとんど確実に」は確率 $1$ の事象上で成り立つことを意味する．
時間集合 $I$ に添字づけられた確率変数族
$(X_t)_{t\in I}$ を確率過程という．
各 $\omega$ に対し $t\mapsto X_t(\omega)$ を標本路という．
確率過程 $X,Y$ が\emph{修正}の関係にあるとは，
各固定した $t$ に対して $P(X_t=Y_t)=1$ が成立することをいう．
$P(X_t=Y_t\text{ for all }t\in I)=1$ のとき，
両過程は\emph{区別不能}であるという．
前者から後者が自動的に従うわけではない．

\begin{definition}
実確率過程 $(X_t)$ がGaussian過程であるとは，
任意の有限個の時刻 $t_1,\dots,t_n$ に対して
$(X_{t_1},\dots,X_{t_n})$ が多変量正規分布に従うことをいう．
退化した正規分布も許す．
\end{definition}

\begin{lemma}[Gaussian系の独立性]\label{lem:gauss-ind}
有限個の同時Gaussianな確率変数
$X_1,\dots,X_m,Y_1,\dots,Y_n$ が
$\operatorname{Cov}(X_i,Y_j)=0$ を全ての $i,j$ で満たすなら，
ベクトル $(X_1,\dots,X_m)$ と
$(Y_1,\dots,Y_n)$ は独立である．
\end{lemma}
\begin{proof}
同時Gaussianベクトルの特性関数は，平均ベクトル $m$
と共分散行列 $\Sigma$ により
\[
 E[e^{i\langle u,Z\rangle}]
 =\exp\!\left(i\langle u,m\rangle
 -\frac12\langle u,\Sigma u\rangle\right)
\]
と表される．
仮定により共分散行列の交差ブロックは零であるから，
$(X,Y)$ の特性関数は $X$ の特性関数と $Y$ の特性関数の積になる．
特性関数の一意性から両ベクトルの独立性を得る．
\end{proof}

\section{Brown運動の定義}

\begin{definition}
実数値過程 $(B_t)_{t\ge0}$ が
\emph{標準Brown運動}であるとは，
次を満たすことをいう．
\begin{enumerate}
\item $B_0=0$ がほとんど確実に成立する．
\item 任意の $0\le t_0<t_1<\cdots<t_n$ に対して
$B_{t_1}-B_{t_0},\dots,B_{t_n}-B_{t_{n-1}}$
は互いに独立である．
\item $0\le s<t$ に対し $B_t-B_s$ は
平均 $0$，分散 $t-s$ の正規分布に従う．
\item 標本路 $t\mapsto B_t(\omega)$ は，
確率 $1$ で連続である．
\end{enumerate}
\end{definition}

\begin{proposition}\label{prop:covariance}
標準Brown運動は中心Gaussian過程であり，
\[
 E[B_sB_t]=\min\{s,t\}
 \qquad(s,t\ge0)
\]
を満たす．逆に，連続な標本路を持つ中心Gaussian過程で
$E[B_sB_t]=\min\{s,t\}$ を満たすものは
標準Brown運動である．
\end{proposition}
\begin{proof}
Brown運動の任意の有限個の値は，
独立な正規増分の線形結合なので同時Gaussianである．
$s\le t$ の場合，$B_t=B_s+(B_t-B_s)$ と
独立増分により
$E[B_sB_t]=E[B_s^2]=s$ となる．

逆に共分散が $\min\{s,t\}$ である中心Gaussian過程を考える．
$E[B_0^2]=0$ なので $B_0=0$ は概至る所成立する．
\[
 E[(B_t-B_s)^2]=t+s-2\min\{s,t\}=t-s
\]
である．互いに交わらない区間に対応する増分の共分散は，
$\min\{s,t\}$ を代入して展開すると $0$ となる．
補題\ref{lem:gauss-ind}から独立性が従う．
従って標準Brown運動の全条件を満たす．
\end{proof}

\section{連続な修正を与える定理}

\begin{theorem}[Kolmogorovの連続性定理]\label{thm:kol-cont}
$T>0$ とし，$X=(X_t)_{0\le t\le T}$ を実確率過程とする．
ある $\alpha,\beta,C>0$ が存在して
\[
 E|X_t-X_s|^\alpha\le C|t-s|^{1+\beta}
 \qquad(0\le s,t\le T)
 \tag{$*$}
\]
が成り立つとする．
このとき $X$ は連続な標本路を持つ修正を持つ．
さらに任意の $\gamma\in(0,\beta/\alpha)$ に対して，
その修正の標本路はほとんど確実に
指数 $\gamma$ のH\"older連続関数となる．
\end{theorem}
\begin{proof}
まず $T=1$ とし，$0<\gamma<\beta/\alpha$ を固定する．
$D_n=\{k2^{-n}:0\le k\le2^n\}$ と置き，
隣接した二点の増分の最大値を
\[
 Z_n=\max_{0\le k<2^n}
 |X_{(k+1)2^{-n}}-X_{k2^{-n}}|
\]
とする．Markovの不等式と $(*)$ により
\[
 P(Z_n>2^{-n\gamma})
 \le 2^n C\,2^{-n(1+\beta)}2^{n\alpha\gamma}
 =C2^{-n(\beta-\alpha\gamma)}.
\]
右辺の $n$ に関する和は収束するから，
Borel--Cantelliの補題により，確率 $1$ で
十分大きい全ての $n$ について
$Z_n\le2^{-n\gamma}$ となる．
有限個の例外的な $n$ を吸収すれば，
ある有限な確率変数 $C_\gamma$ が存在して
\[
 Z_n\le C_\gamma2^{-n\gamma}\qquad(n\ge0)
 \tag{$**$}
\]
を同時に満たす．

$D=\bigcup_nD_n$ と置く．
任意の $u,v\in D$，$0<|u-v|\le2^{-m}$ を取る．
各 $n\ge m$ について
$u_n=2^{-n}\lfloor2^nu\rfloor$，
$v_n=2^{-n}\lfloor2^nv\rfloor$ とする．
$u_m,v_m$ は等しいか隣接しており，
$u_n\to u$，$v_n\to v$ は有限段階で一致する．
隣り合う近似値の差には $(**)$ を使えるから
\begin{align*}
 |X_u-X_v|
 &\le |X_{u_m}-X_{v_m}|\\
 &\quad+\sum_{n=m}^{\infty}
  (|X_{u_{n+1}}-X_{u_n}|
   +|X_{v_{n+1}}-X_{v_n}|)\\
 &\le C_\gamma
 \left(2^{-m\gamma}+2\sum_{n=m+1}^\infty2^{-n\gamma}\right)
 \le C'_\gamma2^{-m\gamma}.
\end{align*}
従って $D$ 上の標本路は一様連続であり，
$[0,1]$ 上の連続関数 $\widetilde X$ に一意に拡張される．
上の評価と $|u-v|$ に対応する $m$ の選択から，
この拡張は指数 $\gamma$ のH\"older連続関数である．

各 $t\in[0,1]$ に対し $t_n\in D$，$t_n\to t$ を取る．
仮定 $(*)$ から $X_{t_n}\to X_t$ は
$L^\alpha$，従って確率収束である．
一方，構成により $X_{t_n}\to\widetilde X_t$
は概至る所成立する．確率極限の一意性から
$\widetilde X_t=X_t$ が概至る所成立する．
よって $\widetilde X$ は $X$ の修正である．

一般の $T>0$ は $t=Tu$ という時間の再尺度化で帰着できる．
また任意の $\gamma<\beta/\alpha$ についての主張は，
$\gamma$ を含む可算稠密な指数列を選び，
それぞれの確率1の集合の共通部分を取ればよい．
\end{proof}

\begin{remark}
証明中の構成は連続な修正を与える．
元の $X$ の標本路が連続であるとは仮定していない．
各 $t$ での概一致と全ての $t$ での概一致は異なる概念である．
\end{remark}

\section{Brown運動の存在}

\begin{theorem}[Brown運動の存在]\label{thm:existence}
ある確率空間上に標準Brown運動が存在する．
\end{theorem}
\begin{proof}
任意の有限個の時刻 $t_1,\dots,t_n\ge0$ に対し，
平均零で共分散行列
$\Sigma_{ij}=\min\{t_i,t_j\}$ を持つGaussianベクトルの
分布を考える．この共分散行列は半正定値である．実際，
任意の実数 $a_1,\dots,a_n$ に対して
\[
 \sum_{i,j=1}^n a_ia_j\min\{t_i,t_j\}
 =\int_0^\infty
   \left(\sum_{i=1}^n a_i1_{[0,t_i]}(u)\right)^2du\ge0.
\]
従ってGaussianベクトルを構成できる．
この有限次元分布族は座標の周辺化と時刻の置換に関して
整合的である．Kolmogorovの拡張定理により，
これらを有限次元分布に持つ確率過程
$(X_t)_{t\ge0}$ が存在する．
命題\ref{prop:covariance}により
$X$ は連続性以外のBrown運動の性質を満たす．

正規確率変数 $Z\sim N(0,\sigma^2)$ について
$E[Z^4]=3\sigma^4$ であるから
\[
 E|X_t-X_s|^4=3|t-s|^2.
\]
任意の $T>0$ について
定理\ref{thm:kol-cont}を $\alpha=4,\beta=1$ で適用し，
$[0,T]$ 上の連続な修正を得る．
整数 $T=1,2,\dots$ ごとに得た修正は，
重なり合う区間の有理数時刻で
元の過程に同時に一致する確率1の集合が存在するため，
連続性によって重なり合う区間上で区別不能である．
従ってそれらを確率1の集合上で貼り合わせると，
$[0,\infty)$ 上で連続な修正 $(B_t)$ を得る．
修正は有限次元分布を変えないので，
$B$ は標準Brown運動である．
\end{proof}

\begin{proposition}[標本路のH\"older性]\label{prop:holder}
標準Brown運動の標本路は，
任意の $T>0$ および $0<\gamma<1/2$ について，
ほとんど確実に $[0,T]$ 上で指数 $\gamma$ の
H\"older連続関数である．
\end{proposition}
\begin{proof}
正規分布の任意の正の偶数次モーメントは有限なので，
偶数 $p\ge4$ に対して
\[
 E|B_t-B_s|^p=C_p|t-s|^{p/2}
\]
が成り立つ．
定理\ref{thm:kol-cont}から指数
$\gamma<(p/2-1)/p=1/2-1/p$ のH\"older性を得る．
$p$ を十分大きく取れば任意の $\gamma<1/2$ を得る．
時間区間については正整数 $T$ に関する可算共通部分を
取ればよい．
\end{proof}

\section{Brown運動の基本変換}

\begin{proposition}\label{prop:scaling}
$B$ を標準Brown運動とする．
\begin{enumerate}
\item 任意の $a>0$ について
$W_t=a^{-1/2}B_{at}$ は標準Brown運動である．
\item 任意の固定した $s\ge0$ について
$W_t=B_{s+t}-B_s$ は標準Brown運動である．
\item $W_t=-B_t$ は標準Brown運動である．
\end{enumerate}
\end{proposition}
\begin{proof}
いずれも標本路の連続性と始点 $0$ は明らかである．
(1) の増分は独立で，分散は
$a^{-1}\operatorname{Var}(B_{at}-B_{au})=t-u$ である．
(2) は定常独立増分の定義から従う．
(3) は中心正規分布の対称性から従う．
\end{proof}

\begin{proposition}[Markov性]\label{prop:markov}
$B$ を標準Brown運動とし，
$\mathcal F_t=\sigma(B_s:0\le s\le t)$
を通常の完備化のもとで考える．
任意の $t,h\ge0$ と有界Borel関数 $f$ に対し
\[
 E[f(B_{t+h})\mid\mathcal F_t]
 =\int_{\mathbb R}f(B_t+z)
   \frac{e^{-z^2/(2h)}}{\sqrt{2\pi h}}\,dz
 \quad(h>0)
\]
が成り立つ．
$h=0$ のときは右辺を $f(B_t)$ と解釈する．
\end{proposition}
\begin{proof}
Brown運動の独立増分により
$B_{t+h}-B_t$ は $\mathcal F_t$ と独立で
$N(0,h)$ に従う．
この独立性は，まず有限個の $B_s$（$s\le t$）
から生成される代数について成り立ち，
$\pi$-$\lambda$ 定理で $\mathcal F_t$ 全体へ拡張できる．
有界Borel関数と独立性を用いる条件付き期待値の公式から
結論を得る．
\end{proof}

\section{停止時刻と強Markov性}

\begin{definition}
$(\mathcal F_t)_{t\ge0}$ をフィルトレーションとする．
$[0,\infty]$ 値確率変数 $\tau$ が停止時刻であるとは，
$\{\tau\le t\}\in\mathcal F_t$ がすべての $t\ge0$ で成り立つことをいう．
$\mathcal F_\tau
=\{A\in\mathcal F:A\cap\{\tau\le t\}\in\mathcal F_t
 \ (t\ge0)\}$ と置く．
\end{definition}

\begin{theorem}[Brown運動の強Markov性]\label{thm:strong-markov}
$B$ を，通常の条件を満たし増分が過去と独立である
フィルトレーションに関する標準Brown運動とする．
$\tau$ をほとんど確実に有限な停止時刻とする．
このとき
\[
 W_t=B_{\tau+t}-B_\tau\qquad(t\ge0)
\]
は標準Brown運動であり，
$\mathcal F_\tau$ と独立である．
\end{theorem}
\begin{proof}
まず $\tau$ が有界とする．
$n\ge1$ に対して
$\tau_n=2^{-n}\lceil2^n\tau\rceil$ と置く．
これは有限個の値を取る停止時刻で，
$\tau\le\tau_n<\tau+2^{-n}$ を満たす．
任意の時刻 $0\le t_1<\cdots<t_m$，
有界連続関数 $g:\mathbb R^m\to\mathbb R$，
$A\in\mathcal F_\tau$ に対して
\[
 E\!\left[1_Ag(B_{\tau_n+t_1}-B_{\tau_n},
 \dots,B_{\tau_n+t_m}-B_{\tau_n})\right]
 =P(A)E[g(B_{t_1},\dots,B_{t_m})].
 \tag{$*$}
\]
実際，$\tau_n$ の取り得る各値 $s$ について
$A\cap\{\tau_n=s\}\in\mathcal F_s$ であり，
$s$ 以降の増分は $\mathcal F_s$ と独立で，
定常増分性から右辺の正規分布と一致する．
各値について足し合わせれば $(*)$ を得る．

標本路の連続性から $\tau_n\to\tau$ とともに
各 $B_{\tau_n+t_j}-B_{\tau_n}$ は
$B_{\tau+t_j}-B_\tau$ に概収束する．
有界収束定理で $(*)$ の左辺に極限を取れば，
$W$ の有限次元分布が $B$ のそれに一致し，
かつその全ての有限次元分布が
$\mathcal F_\tau$ と独立であることを得る．
$W$ は連続標本路を持つから標準Brown運動であり，
有理時刻の評価写像が生成する $\sigma$-加法族を通じて
過程全体が $\mathcal F_\tau$ と独立である．

一般の $\tau<\infty$ については
$\tau\wedge N$ に以上を適用し，
$A\cap\{\tau\le N\}\in\mathcal F_{\tau\wedge N}$
（$A\in\mathcal F_\tau$）を用いる．
$N\to\infty$ として有界収束定理を適用すれば
同じ有限次元分布の独立性を得る．
\end{proof}

\section{反射原理と到達時刻}

\begin{definition}
$a>0$ とし
$\tau_a=\inf\{t\ge0:B_t=a\}$ と置く．
該当する時刻がなければ $\tau_a=\infty$ とする．
標本路の連続性から
$\{\tau_a\le t\}=\{\max_{0\le s\le t}B_s\ge a\}$
であり，$\tau_a$ は停止時刻である．
\end{definition}

\begin{theorem}[反射原理]\label{thm:reflection}
$a>0$，$t>0$ とすると
\[
 P\left(\max_{0\le s\le t}B_s\ge a\right)
 =2P(B_t\ge a)
 =2\left(1-\Phi\!\left(\frac a{\sqrt t}\right)\right),
\]
ただし
$\Phi(u)=(2\pi)^{-1/2}\int_{-\infty}^u e^{-v^2/2}\,dv$
とする．
\end{theorem}
\begin{proof}
$A=\{\tau_a\le t\}$ と置く．
$B_t\ge a$ なら必ず $A$ が生じるため
\[
 P(A)=P(B_t\ge a)+P(A,B_t<a).
\]
強Markov性とBrown運動の対称性により，
$\tau_a$ 以降の標本路を $a$ で反射した過程
\[
 \widetilde B_s=
 \begin{cases}
 B_s&(s\le\tau_a),\\
 2a-B_s&(s>\tau_a)
 \end{cases}
\]
は $B$ と同じ法則を持つ．
これは停止時刻以降の独立なBrown運動の増分を
符号反転しても法則が変わらないことによる．
$A$ 上では
$B_t<a$ と $\widetilde B_t>a$ は同値である．
反射写像は反転操作に関して逆写像自身なので，
\[
 P(A,B_t<a)=P(A,B_t>a).
\]
$B_t>a$ なら $A$ が生じることから
$P(A,B_t>a)=P(B_t>a)$ である．
$B_t$ は連続分布を持つので
$P(B_t=a)=0$ を用いて
$P(A)=2P(B_t\ge a)$ を得る．
最後に $B_t\sim N(0,t)$ を用いる．
\end{proof}

\begin{corollary}[到達時刻の分布]\label{cor:hitting}
$a>0$ に対して $\tau_a<\infty$ が概至る所で成立する．
さらに $\tau_a$ は $(0,\infty)$ 上の密度
\[
 f_{\tau_a}(t)
 =\frac{a}{\sqrt{2\pi}\,t^{3/2}}
   \exp\!\left(-\frac{a^2}{2t}\right)
 \qquad(t>0)
\]
を持つ．
\end{corollary}
\begin{proof}
定理\ref{thm:reflection}により
$P(\tau_a\le t)=2(1-\Phi(a/\sqrt t))$ である．
$t\to\infty$ で右辺は $1$ に収束するので
$\tau_a<\infty$ が概至る所である．
$t>0$ について分布関数を微分すれば表示の密度を得る．
\end{proof}

\section{Brown運動の二次変分}

\begin{theorem}[二次変分]\label{thm:qv}
$T>0$ とし，
$\pi_n=\{0=t_0^{(n)}<\cdots<t_{k_n}^{(n)}=T\}$
を分割列，$|\pi_n|=\max_j(t_{j+1}^{(n)}-t_j^{(n)})$
を最大分割幅とする．
$|\pi_n|\to0$ なら
\[
 Q_n=\sum_{j=0}^{k_n-1}
 (B_{t_{j+1}^{(n)}}-B_{t_j^{(n)}})^2
 \longrightarrow T
\]
が $L^2$ で成立する．
特に等分割で $k_n=2^n$ の場合，
収束はほとんど確実にも成立する．
\end{theorem}
\begin{proof}
分割の幅を $\Delta_j=t_{j+1}^{(n)}-t_j^{(n)}$ と置く．
Brown運動の独立増分と正規分布の四次モーメントより
\[
 E[Q_n]=\sum_j\Delta_j=T,\qquad
 \operatorname{Var}(Q_n)
 =\sum_j2\Delta_j^2\le2T|\pi_n|.
\]
従って
$E|Q_n-T|^2\le2T|\pi_n|\to0$ である．
等分割 $|\pi_n|=T2^{-n}$ では
$\sum_nE|Q_n-T|^2<\infty$ だから，
任意の $\varepsilon>0$ に対する
Chebyshevの不等式とBorel--Cantelliの補題により
$Q_n\to T$ がほとんど確実に成立する．
\end{proof}

\begin{corollary}[Brown運動の変動]\label{cor:unbounded-variation}
$T>0$ について，標準Brown運動の標本路は
$[0,T]$ 上でほとんど確実に有界変動でない．
\end{corollary}
\begin{proof}
もし連続関数 $f:[0,T]\to\mathbb R$ の全変動
$V_T(f)$ が有限なら，任意の分割 $\pi$ に対して
\[
 \sum_{[u,v]\in\pi}|f(v)-f(u)|^2
 \le\left(\max_{[u,v]\in\pi}|f(v)-f(u)|\right)V_T(f).
\]
分割幅が零へ向かうと，$f$ の一様連続性から右辺は $0$ へ向かう．
一方，定理\ref{thm:qv}によりBrown運動の等分割の二次変分は
概至る所 $T>0$ に収束する．従って有界変動ではあり得ない．
\end{proof}

\begin{proposition}[各固定時刻での非微分可能性]
$t\ge0$ を固定する．標準Brown運動の標本路が
$t$ で有限な右微分係数を持つ確率は $0$ である．
\end{proposition}
\begin{proof}
$h>0$ に対して
\[
 \frac{B_{t+h}-B_t}{h}\sim N(0,h^{-1})
\]
である．
任意の $M>0$ について
\[
 P\left(\left|\frac{B_{t+h}-B_t}{h}\right|\le M\right)
 =P(|N(0,1)|\le M\sqrt h)\longrightarrow0.
\]
もし有限な右微分係数が存在するなら，
例えば $h_n=1/n$ に沿う差商は有限値に収束するため，
ある整数 $M$ について十分大きい $n$ で
差商の絶対値が $M$ 以下となる．
これは上の確率評価とFatouの補題を適用した
事象の極限に反する．
従って有限な右微分係数が存在する確率は零である．
\end{proof}

\begin{remark}
最後の命題は，各固定時刻で微分可能である確率が零であることを
示している．「ほとんど確実に全ての時刻で微分不可能」という
さらに強い命題をこれだけから導くことはできない．
\end{remark}

\section{Wiener測度}

\begin{proposition}\label{prop:wiener}
$T>0$ とし，
$C_0([0,T])=\{f\in C([0,T],\mathbb R):f(0)=0\}$
に一様ノルムによるBorel $\sigma$-加法族を入れる．
Brown運動の連続標本路を値に持つ写像
$\omega\mapsto B_\cdot(\omega)$ は
$C_0([0,T])$ 値確率変数である．
その法則をWiener測度といい，
Brown運動の有限次元分布によって一意に決まる．
\end{proposition}
\begin{proof}
$C_0([0,T])$ の一様ノルムによる開球
$B(f,r)$ は
\[
 \{g:\norm{g-f}_\infty<r\}
 =\bigcup_{m\ge1}\bigcap_{q\in\mathbb Q\cap[0,T]}
  \{g:|g(q)-f(q)|\le r-1/m\}
\]
という形で表せる（$r-1/m>0$ の項だけを使う）．
従って有理数時刻での評価写像が生成する
$\sigma$-加法族はBorel $\sigma$-加法族と一致する．
さらに $C_0([0,T])$ は可分である．
実際，有理数時刻（および $T$）を分点とし，
各分点の値を有理数とする折れ線関数で，
始点の値を $0$ とするもの全体は可算集合である．
連続関数の一様連続性と折れ線補間を用いれば，
この集合は $C_0([0,T])$ で稠密である．
従って可算な開球の基が存在し，
任意の開集合は高々可算個の基本開球の和集合として表せる．
各 $B_q$ は確率変数であるから，
標本路写像による全ての基本開球の逆像が可測であり，
任意の開集合の逆像も可測となる．
よって標本路写像は可測である．
また有限次元分布が与えられれば，
有理数時刻の座標に関する円柱集合の確率が決まる．
これらの円柱集合は上のBorel $\sigma$-加法族を生成するため，
確率測度の一意性定理からWiener測度は一意である．
\end{proof}

\begin{remark}
標本路が連続であるという条件を持つ確率過程は，
有限次元分布だけでなく関数空間上の確率測度として扱える．
この立場は確率微分方程式や確率積分の構成における
基本的な枠組みを与える．
\end{remark}

\section{到達課題の完全解答}
\noindent\textbf{Brown運動の共分散計算．}
\begin{proposition}
標準Brown運動 $(B_t)_{t\ge0}$ に対し
\[
E[B_sB_t]=\min\{s,t\},\qquad
E[(B_t-B_s)^2]=|t-s|
\]
が成り立つ．
\end{proposition}
\begin{proof}
$s\le t$ としてよい．
$B_t=B_s+(B_t-B_s)$ であり，
増分は $B_s$ と独立で平均零だから
\[
E[B_sB_t]
=E[B_s^2]+E[B_s]E[B_t-B_s]
=s.
\]
増分 $B_t-B_s$ は平均零，分散 $t-s$
の正規分布なので二乗平均は $t-s$ である．
対称性から一般の場合も従う．
\end{proof}
\noindent\textbf{到達確認．}
$E[(B_3-B_1)B_1]$ を求めよ．
\par\noindent\textbf{解答．}
増分 $B_3-B_1$ と $B_1$ は独立でそれぞれ
平均零なので期待値は0である．
\end{document}
